
\documentclass[11pt]{article}

\input{../../tex/defs.tex}

% Useful syntax commands
\newcommand{\jarr}[1]{\left[#1\right]}   % \jarr{x: y} = {x: y}
\newcommand{\jobj}[1]{\left\{#1\right\}} % \jobj{1, 2} = [1, 2]
\newcommand{\pgt}[1]{\, > {#1}}          % \pgt{1} = > 1
\newcommand{\plt}[1]{\, < {#1}}          % \plt{2} = < 2
\newcommand{\peq}[1]{\, = {#1}}          % \peq{3} = = 3
\newcommand{\prop}[1]{\langle{#1}\rangle}% \prop{x} = <x>
\newcommand{\matches}[2]{{#1}\sim{#2}}   % \matches{a}{b} = a ~ b
\newcommand{\aeps}{\varepsilon}          % \apes = epsilon
\newcommand{\akey}[2]{.{#1}\,{#2}}       % \akey{s}{a} = .s a
\newcommand{\aidx}[2]{[#1]\,{#2}}        % \aidx{i}{a} = [i] a
\newcommand{\apipe}[1]{\mid {#1}}        % \apipe{a} = | a

% Solution-specific helpers
\newcommand{\amap}[2]{[\,{#1}\,]\,{#2}}          % map accessor: [a] a'
\newcommand{\tup}[1]{\langle {#1} \rangle}        % tuple schema <t0,...,tk>
\newcommand{\sat}[2]{{#1} \vDash {#2}}            % number satisfies property: n |= p
\newcommand{\valid}[3]{{#1} \vdash {#2} : {#3}}   % validity w/ result schema: t |- a : t'

% Other useful syntax commands:
%
% \msf{x} = x (not italicised)
% \falset = false
% \truet = true
% \tnum = num
% \tbool = bool
% \tstr = str

\newtheorem{theorem}{Theorem}
\newtheorem{lemma}{Lemma}

\begin{document}

\hwtitle
  {Assignment 1: JSafe}
  {appleweiping} %% REPLACE THIS WITH YOUR NAME/ID

We formalise a small typed query language over JSON. A \emph{schema} $\tau$
describes a set of JSON documents; an \emph{accessor} $a$ is a path/query.
The goal (Problem 2) is to prove that a well-typed accessor, run on any
document matching the schema, always terminates with a result.

Throughout, $n,m$ range over (integer) numbers, $s$ over strings, $j$ over JSON
values, and $\str{X}$ denotes zero-or-more repetitions of $X$. Objects are
finite maps from distinct string keys to values; array/object element order is
significant only where stated.

\problem{Problem 1: Schemas}

\subsection*{Part 1: Syntax}

\begin{alignat*}{2}
\msf{JSON}~j ::= \qamp \truet \mid \falset \mid n \mid s
  \mid \jarr{j_0, \ldots, j_{k-1}}
  \mid \jobj{s_0 : j_0, \ldots, s_{k-1} : j_{k-1}} \\[0.4em]
\msf{Property}~p ::= \qamp \aeps
  \mid (\pgt{n})\,p
  \mid (\plt{n})\,p
  \mid (\peq{n})\,p \\[0.4em]
\msf{Schema}~\tau ::= \qamp \tbool
  \mid \tnum \prop{p}
  \mid \tstr
  \mid \jarr{\tau}
  \mid \tup{\tau_0, \ldots, \tau_{k-1}}
  \mid \jobj{s_0 : \tau_0, \ldots, s_{k-1} : \tau_{k-1}}
\end{alignat*}

A property $p$ is a conjunction of numeric constraints. A number schema
$\tnum\prop{p}$ matches numbers satisfying \emph{all} of them.  There are two
array schemas: the \emph{homogeneous} array $\jarr{\tau}$ (any length, every
element matches $\tau$) and the \emph{tuple} $\tup{\tau_0,\ldots,\tau_{k-1}}$
(exactly $k$ elements, the $i$-th matching $\tau_i$).  An object schema fixes a
set of keys and the schema of each.

\subsection*{Part 2: Matching}

We first define when a number satisfies a property, $\sat{n}{p}$
(``$n$ satisfies every constraint in $p$''):

\begin{mathpar}
\ir{P-Emp}{\ }{\sat{n}{\aeps}}

\ir{P-Gt}{n > m \\ \sat{n}{p}}{\sat{n}{(\pgt{m})\,p}}

\ir{P-Lt}{n < m \\ \sat{n}{p}}{\sat{n}{(\plt{m})\,p}}

\ir{P-Eq}{n = m \\ \sat{n}{p}}{\sat{n}{(\peq{m})\,p}}
\end{mathpar}

The matching judgment $\matches{j}{\tau}$ reads ``JSON value $j$ adheres to
schema $\tau$'':

\begin{mathpar}
\ir{S-Bool-False}{\ }{\matches{\falset}{\tbool}}

\ir{S-Bool-True}{\ }{\matches{\truet}{\tbool}}

\ir{S-Num}{\sat{n}{p}}{\matches{n}{\tnum\prop{p}}}

\ir{S-Str}{\ }{\matches{s}{\tstr}}

\ir{S-Arr}
  {\matches{j_i}{\tau} \quad (\text{for all } 0 \le i < k)}
  {\matches{\jarr{j_0, \ldots, j_{k-1}}}{\jarr{\tau}}}

\ir{S-Tup}
  {\matches{j_i}{\tau_i} \quad (\text{for all } 0 \le i < k)}
  {\matches{\jarr{j_0, \ldots, j_{k-1}}}{\tup{\tau_0, \ldots, \tau_{k-1}}}}

\ir{S-Obj}
  {\matches{j_i}{\tau_i} \quad (\text{for all } 0 \le i < k)}
  {\matches{\jobj{s_0 : j_0, \ldots, s_{k-1} : j_{k-1}}}
           {\jobj{s_0 : \tau_0, \ldots, s_{k-1} : \tau_{k-1}}}}
\end{mathpar}

(S-Arr allows any length; S-Tup requires the value to have exactly the $k$
keys/indices of the tuple; S-Obj requires the same key set, matched pointwise.)

\problem{Problem 2: Accessors}

\subsection*{Part 1: Operational semantics}

An accessor is a path that consumes one document and produces one document:

\begin{alignat*}{1}
\msf{Accessor}~a ::= \qamp \aeps
  \mid \akey{s}{a}
  \mid \aidx{n}{a}
  \mid \amap{a}{a}
\end{alignat*}

$\aeps$ is the identity (final) accessor; $\akey{s}{a}$ looks up key $s$ then
continues with $a$; $\aidx{n}{a}$ indexes an array then continues; and
$\amap{a_1}{a_2}$ (the bracket contains an \emph{accessor} $a_1$, not a numeral)
maps $a_1$ over every element of an array and continues with $a_2$.

A configuration is a pair $(a, j)$. The small-step relation
$\steps{(a,j)}{(a',j')}$ (and its reflexive--transitive closure
$\evals{(a,j)}{(a',j')}$) is:

\begin{mathpar}
\ir{E-Key}
  {\ }
  {\steps
    {(\akey{s_i}{a},\ \jobj{s_0 : j_0, \ldots, s_{k-1} : j_{k-1}})}
    {(a,\ j_i)}}

\ir{E-Idx}
  {0 \le n < k}
  {\steps
    {(\aidx{n}{a},\ \jarr{j_0, \ldots, j_{k-1}})}
    {(a,\ j_n)}}

\ir{E-Map}
  {\evals{(a_1,\ j_i)}{(\aeps,\ j_i')} \quad (\text{for all } 0 \le i < k)}
  {\steps
    {(\amap{a_1}{a_2},\ \jarr{j_0, \ldots, j_{k-1}})}
    {(a_2,\ \jarr{j_0', \ldots, j_{k-1}'})}}
\end{mathpar}

A configuration $(\aeps, j)$ is \emph{final}: no rule applies, and it is the
intended answer. E-Key fires only when $s_i$ is actually a key of the object;
E-Idx only when the index is in bounds; E-Map runs the inner accessor $a_1$ to
completion on each element (its premises are themselves multi-step evaluations)
and rebuilds the array. (One may equivalently split E-Map into an empty-array
base case and a one-element inductive step; we use the multi-step form for
brevity.)

\subsection*{Part 2: Validity, and the Accessor Safety theorem}

Not every accessor makes sense on every schema: $\akey{s}{a}$ needs $s$ to be a
key, $\aidx{n}{a}$ needs the array to have $> n$ elements, and $\amap{a_1}{a_2}$
needs $a_1$ to be valid on the element schema. We therefore give a
\emph{validity} judgment that also records the schema of the result,
$\valid{\tau}{a}{\tau'}$, read ``accessor $a$ is valid on a document of schema
$\tau$ and yields a document of schema $\tau'$'':

\begin{mathpar}
\ir{V-Eps}
  {\ }
  {\valid{\tau}{\aeps}{\tau}}

\ir{V-Key}
  {\valid{\tau_i}{a}{\tau'}}
  {\valid{\jobj{s_0 : \tau_0, \ldots, s_{k-1} : \tau_{k-1}}}{\akey{s_i}{a}}{\tau'}}

\ir{V-Idx}
  {0 \le n < k \\ \valid{\tau_n}{a}{\tau'}}
  {\valid{\tup{\tau_0, \ldots, \tau_{k-1}}}{\aidx{n}{a}}{\tau'}}

\ir{V-Map}
  {\valid{\tau}{a_1}{\tau_1} \\ \valid{\jarr{\tau_1}}{a_2}{\tau'}}
  {\valid{\jarr{\tau}}{\amap{a_1}{a_2}}{\tau'}}
\end{mathpar}

We write $\matches{a}{\tau}$ (``$a$ is valid on $\tau$'') to mean that
$\valid{\tau}{a}{\tau'}$ holds for some result schema $\tau'$.

\begin{theorem}[Accessor Safety]
For all $a$, $j$, $\tau$: if $\matches{a}{\tau}$ and $\matches{j}{\tau}$, then
there exists $j'$ such that $\evals{(a, j)}{(\aeps, j')}$.
\end{theorem}

We prove a slightly stronger statement, which additionally tracks the schema of
the output. Accessor Safety is the special case obtained by forgetting the last
conjunct.

\begin{lemma}[Safety with preservation]
\label{lem:safe}
If $\valid{\tau}{a}{\tau'}$ and $\matches{j}{\tau}$, then there exists $j'$ such
that $\evals{(a, j)}{(\aeps, j')}$ and $\matches{j'}{\tau'}$.
\end{lemma}

\begin{proof}
By (structural) rule induction on the derivation $\mathcal{D}$ of
$\valid{\tau}{a}{\tau'}$. In each case we do case analysis on the last rule of
$\mathcal{D}$; the inductive hypothesis (IH) may be applied to any
sub-derivation of $\mathcal{D}$, whose accessor is a strict subterm of $a$.

\textbf{Case V-Eps.} Then $a = \aeps$ and $\tau' = \tau$. Take $j' = j$: we have
$\evals{(\aeps, j)}{(\aeps, j)}$ in zero steps, and $\matches{j}{\tau'}$ holds by
assumption.

\textbf{Case V-Key.} Then $a = \akey{s_i}{a_0}$,
$\tau = \jobj{s_0 : \tau_0, \ldots, s_{k-1}:\tau_{k-1}}$, and the premise is
$\valid{\tau_i}{a_0}{\tau'}$. Since $\matches{j}{\tau}$, inversion on the
matching rules forces the last rule to be S-Obj, so
$j = \jobj{s_0 : j_0, \ldots, s_{k-1} : j_{k-1}}$ with the same keys and
$\matches{j_l}{\tau_l}$ for every $l$; in particular $\matches{j_i}{\tau_i}$.
By E-Key, $\steps{(\akey{s_i}{a_0}, j)}{(a_0, j_i)}$. By the IH applied to
$\valid{\tau_i}{a_0}{\tau'}$ with $\matches{j_i}{\tau_i}$, there is $j'$ with
$\evals{(a_0, j_i)}{(\aeps, j')}$ and $\matches{j'}{\tau'}$. Prepending the
single E-Key step gives $\evals{(a, j)}{(\aeps, j')}$.

\textbf{Case V-Idx.} Then $a = \aidx{n}{a_0}$,
$\tau = \tup{\tau_0, \ldots, \tau_{k-1}}$, and the premises are $0 \le n < k$ and
$\valid{\tau_n}{a_0}{\tau'}$. By inversion of $\matches{j}{\tau}$ (rule S-Tup),
$j = \jarr{j_0, \ldots, j_{k-1}}$ with exactly $k$ elements and
$\matches{j_l}{\tau_l}$ for all $l$; in particular $\matches{j_n}{\tau_n}$ and,
since $0 \le n < k$, index $n$ is in bounds. By E-Idx,
$\steps{(\aidx{n}{a_0}, j)}{(a_0, j_n)}$. By the IH on
$\valid{\tau_n}{a_0}{\tau'}$ with $\matches{j_n}{\tau_n}$, there is $j'$ with
$\evals{(a_0, j_n)}{(\aeps, j')}$ and $\matches{j'}{\tau'}$. Prepend the E-Idx
step.

\textbf{Case V-Map.} Then $a = \amap{a_1}{a_2}$, $\tau = \jarr{\tau_0}$ (writing
the element schema as $\tau_0$), and the premises are
$\valid{\tau_0}{a_1}{\tau_1}$ and $\valid{\jarr{\tau_1}}{a_2}{\tau'}$. By
inversion of $\matches{j}{\tau}$ (rule S-Arr), $j = \jarr{j_0, \ldots, j_{k-1}}$
for some $k \ge 0$ with $\matches{j_i}{\tau_0}$ for every $i$.

For each $i$, apply the IH to the sub-derivation $\valid{\tau_0}{a_1}{\tau_1}$
(whose accessor $a_1$ is a strict subterm of $a$) with $\matches{j_i}{\tau_0}$:
we obtain $j_i'$ with $\evals{(a_1, j_i)}{(\aeps, j_i')}$ and
$\matches{j_i'}{\tau_1}$. These are exactly the premises of E-Map, so
\[
  \steps{(\amap{a_1}{a_2}, \jarr{j_0, \ldots, j_{k-1}})}
        {(a_2, \jarr{j_0', \ldots, j_{k-1}'})}.
\]
Moreover $\matches{\jarr{j_0', \ldots, j_{k-1}'}}{\jarr{\tau_1}}$ by S-Arr, since
every $\matches{j_i'}{\tau_1}$. Now apply the IH to the sub-derivation
$\valid{\jarr{\tau_1}}{a_2}{\tau'}$ (its accessor $a_2$ is a strict subterm of
$a$) with this matching: there is $j'$ with
$\evals{(a_2, \jarr{j_0', \ldots, j_{k-1}'})}{(\aeps, j')}$ and
$\matches{j'}{\tau'}$. Prepending the single E-Map step yields
$\evals{(a, j)}{(\aeps, j')}$.

In every case we produced the required $j'$ with $\matches{j'}{\tau'}$, closing
the induction.
\end{proof}

\begin{proof}[Proof of Accessor Safety]
Assume $\matches{a}{\tau}$ and $\matches{j}{\tau}$. By definition of
$\matches{a}{\tau}$ there is some $\tau'$ with $\valid{\tau}{a}{\tau'}$. Applying
Lemma~\ref{lem:safe} yields $j'$ with $\evals{(a, j)}{(\aeps, j')}$ (and, as a
bonus, $\matches{j'}{\tau'}$), which is exactly what the theorem asserts.
\end{proof}

\end{document}
